\section{シミュレーション結果}
\label{sec:simulation_results}

本節では，前節で述べたシミュレーション環境および軟弱地面モデルを用いて実施した歩行シミュレーションの結果について述べる．
従来手法と提案手法を同一条件下で比較し，床沈み込みが生じる状況における歩行安定性の変化を評価する．

\subsection{実験条件}
\label{subsec:sim_conditions}

本シミュレーションでは，歩行周期，歩幅，支持脚切り替え条件などを固定し，床モデルおよび制御手法のみを変更して実験を行った．
シミュレーション条件の詳細を Table~\ref{tab:sim_conditions} に示す．

% （ここに実験条件表）

\subsection{評価指標}
\label{subsec:sim_metrics}

歩行安定性を評価するため，以下の指標を用いた．
\begin{itemize}
  \item ZMP 軌道および支持多角形との関係
  \item 重心（CoM）の位置および高さ変動
  \item 歩行中の転倒の有無
\end{itemize}
これらの指標を用いて，床沈み込みが歩行挙動に与える影響を定量的および定性的に評価した．

\subsection{従来手法による歩行結果}
\label{subsec:sim_conventional}

まず，床沈み込みを考慮しない従来手法を用いた場合の歩行結果を示す．
軟弱地面上においては，支持脚接地時に床沈み込みが発生し，ZMP の変動が大きくなる傾向が確認された．
特に片脚支持期において，ZMP が支持多角形境界付近に移動する挙動が見られた．

% （ZMP・CoM の時系列グラフ）

\subsection{提案手法による歩行結果}
\label{subsec:sim_proposed}

次に，床沈み込みを考慮する提案手法を適用した場合の歩行結果を示す．
提案手法では，床沈み込みに応じて CoM 高さ参照が調整され，ZMP の変動が抑制される傾向が確認された．
その結果，片脚支持期においても安定した歩行が維持された．

% （ZMP・CoM の時系列グラフ）

\subsection{床パラメータの違いによる影響}
\label{subsec:sim_param_effect}

異なる床パラメータを用いたシミュレーション結果を比較することで，
床沈み込み特性の違いが歩行挙動に与える影響を評価した．
ばね定数およびダンパー定数の違いにより，CoM 高さ変動および ZMP 応答に差が生じることが確認された．

\subsection{小結}
\label{subsec:sim_summary}

以上のシミュレーション結果より，床沈み込みを考慮した提案手法は，
軟弱地面上において ZMP 変動を抑制し，歩行安定性を向上させることが示された．
次章では，これらの結果を実機実験により検証する．
